\documentclass[12pt]{article}

\begin{document}

%Capa
%Contracapa

section : {Introdução}
A cidade de Florapolis é conhecida por sua rica cultura, belas paisagens e diversidade de atrações turísticas.  
Neste documento exploraremos diferentes aspectos da cidade, incluindo história, gastronomia, pontos turísticos e análise de crescimento populacional (ver Seção ref:historia e Seção ref:matematica).

section: História
Florapolis foi fundada no século XIX e possui uma arquitetura que combina estilos colonial e moderno.  
O desenvolvimento urbano ao longo dos anos pode ser estudado através de dados populacionais e projeções matemáticas.

subsection: Marcos históricos
Alguns pontos históricos incluem:
     Museu da Cidade
     Praça Central
     Teatro Municipal
     Biblioteca Pública


subsection: Curiosidades matemáticas
O crescimento populacional pode ser estimado pela fórmula exponencial:


P(t) = P0 e{rt} 

onde P0 é a população inicial, r é a taxa de crescimento, e t é o tempo em anos.  
Se considerarmos diferentes bairros com taxas distintas, podemos modelar:

P1(t) = 20000  e{0.015 t} 
P2(t) = 15000  e{0.02 t} 
P3(t) = 10000  e{0.025 t}

A soma dessas populações dá a população total da cidade, P{total}(t) = P1(t) + P2(t) + P3(t), que pode ser útil para planejar serviços públicos (ver Equação ref:bairros).

section: Turismo
As atrações turísticas podem, de acordo com %citação: Knuth, 
ser organizadas por categorias:
     Museus
     Parques e áreas verdes
     Eventos culturais
     Gastronomia típica

subsection: Tabelas
Comparação de pontos turísticos:


Atração  Tipo  Visitantes por ano 
Museu da Cidade  Museu  50.000 
Praça Central  Área Pública  120.000 
Teatro Municipal  Teatro  30.000 
Parque das Flores  Parque  80.000 

Como podemos ver na Tabela ref:atracoes, a Praça Central é o local mais visitado, e o Parque das Flores atrai muitos visitantes durante eventos culturais.

subsection: Figuras
Exemplo de figura (substitua pelo seu arquivo):

% COLOQUE SUA FIGURA AQUI

A Figura ref: florapolis ilustra a beleza natural da cidade, complementando as informações históricas discutidas na Seção ref: historia.

section: Gastronomia
A cidade é famosa por pratos típicos, como bolo de flores e chá de ervas locais.  
Segundo estudos citacão: lamport1994, a gastronomia é um ponto forte do turismo cultural (ver Seção ref: turismo).

section: Análise matemática complementar
Podemos analisar a densidade populacional aproximada de cada bairro usando:


Di(t) = {Pi(t)}{Ai}

onde Ai é a área do bairro i.  
Se considerarmos os três bairros citados na Seção ref: matematica, obtemos:

D1(t) = P1(t)10 
D2(t) = P2(t)8 
D3(t) = P3(t)5 

Essa análise ajuda no planejamento urbano e na logística de transporte, complementando o que foi discutido na Tabela ref: atracoes.

section: Conclusão
Florapolis combina cultura, história, beleza natural e excelente gastronomia.  
O uso de modelos matemáticos (ver Equações ref: populacao, ref: bairros e ref: densidade) permite planejar melhor o crescimento e os serviços da cidade, integrando turismo e desenvolvimento urbano.  
Para mais detalhes sobre atrações turísticas, consulte a Seção ref: turismo e Figura ref: florapolis.

% COLOQUE AQUI AS REFERÊNCIAS.

\end{document}
